La distribución de clientes en mora por sector económico muestra diferencias moderadas entre las categorías analizadas. El sector Comercio registra la mayor cantidad de clientes en mora con
42 casos
, mientras que Industria presenta la menor con
25 casos
, lo que representa una diferencia de
17 clientes
entre ambos extremos. En términos generales, las variaciones observadas son relativamente reducidas considerando el tamaño de la cartera total.
Asimismo, los sectores que concentran un mayor número de clientes también presentan una mayor cantidad absoluta de casos de mora. Este comportamiento sugiere que las diferencias observadas podrían estar asociadas al tamaño de la cartera de clientes en cada sector más que a un efecto propio de la actividad económica. No obstante, con esta visualización no es posible establecer una relación causal o afirmar que el sector económico sea un factor determinante del incumplimiento.
Como resultado, las estrategias de gestión del riesgo crediticio no deberían diseñarse únicamente en función del sector económico. Se recomienda complementar el análisis incorporando variables con mayor capacidad explicativa, como la relación cuota-ingreso, el monto del préstamo, la antigüedad laboral o el historial crediticio, con el fin de identificar los factores que influyen de manera más significativa en la probabilidad de mora.